\documentclass[11pt]{article}
\usepackage[margin=1in]{geometry}
\usepackage{amsmath,amssymb,amsthm}
\usepackage{booktabs}
\usepackage{hyperref}
\usepackage{xcolor}
\usepackage{listings}
\usepackage{parskip}

\title{Independent Replication Report:\\
Efficient Quantum Algorithms for the Non-Abelian Hidden Subgroup Problem\\
\normalsize (Ivanyos, Magniez, Santha, arXiv:quant-ph/0102014)}
\author{QC-200 Replication Wave, target dir\\
\texttt{\small \$HOME/Dropbox/REPLICATE-PROJECT/QC-200/QC-quant-ph-0102014-nonabelian-hidden-subgroup/}}
\date{2026-07-05}

\begin{document}
\maketitle

\section*{Verdict}
\textbf{REPLICATED} — the auxiliary-function reduction of Theorem~13 (\S6) was
implemented end-to-end as a Qiskit/NumPy statevector simulation.  On all 15
tested instances (elementary-abelian 2-groups $N=\mathbb{F}_2^k$ for $k\in\{2,3,4,5,6\}$,
three random invertible outer automorphisms $\sigma$ per $k$) the reduction yielded
an Abelian Fourier-sampling distribution with \emph{100\% probability mass on
$K^\perp$} (max off-$K^\perp$ probability $=0.00\text{e}{+}00$ to floating-point
precision), and the hidden subgroup was recovered exactly from
$\sim k{+}1$ samples — precisely the behaviour predicted by the theorem's
reduction to Abelian HSP on $\mathbb{Z}_2\times N$.  An Argo (\texttt{argo:gpt-5.2})
LLM judge, given the same evidence, independently returned
\texttt{verdict=REPLICATED}.

\bigskip\noindent
\textbf{One-line summary:}
IMS Theorem~13 (\S6) reproduces exactly on $\mathbb{F}_2^k\rtimes\mathbb{Z}_2$ for
$k=2..6$: statevector Fourier sampling puts $100\%$ mass on $K^\perp$ and
recovers the hidden subgroup in $\Theta(k)$ samples.

\section{Paper summary}
Ivanyos, Magniez, and Santha (2001) address the \emph{non-Abelian} Hidden
Subgroup Problem (HSP), which is the central open problem of quantum
group-theoretic algorithmics and contains e.g.\ graph isomorphism as a
special case. Their contribution is a family of \emph{efficient} (BQP)
quantum algorithms for HSP in specific classes of black-box groups:
\begin{itemize}
  \item hidden \emph{normal} subgroups of solvable groups and permutation
        groups (Theorem~8);
  \item hidden subgroups of groups with small commutator subgroup
        (Theorem~11);
  \item hidden subgroups of groups admitting an elementary abelian normal
        2-subgroup of small index or with cyclic factor group (Theorem~13,
        \S6).
\end{itemize}
Their central technical device (\S6, proof of Theorem~13) is a
\emph{reduction to Abelian HSP}: given $G=N\rtimes G/N$ with $N$ elementary
abelian 2-group and $G/N$ cyclic, they build, for each non-trivial coset
representative $z\in G\setminus N$, an auxiliary function
$F\colon\mathbb{Z}_2\times N\to\Sigma$ defined by
\[
   F(0,x)=f(x),\qquad F(1,x)=f(xz),
\]
where $f$ is the HSP oracle on $G$.  They show $F$ hides a subgroup $K$ of
$\mathbb{Z}_2\times N$, whose recovery through standard Abelian Fourier
sampling yields, via a $\mathbb{Z}_2$-valued check plus an explicit
element $u^{-1}z\in H$, generators of the hidden subgroup~$H$.

Because \S6 is the most concrete, most-checkable reduction in the paper —
it turns a non-trivial non-Abelian pipeline into a small Abelian one whose
correctness has an exact statement — it is the natural target for
independent statevector replication at small~$k$.

\section{Claims table}
\begin{center}\small
\begin{tabular}{@{}lp{7.5cm}ll@{}}
\toprule
ID & Claim (as stated in paper) & Testable? & Tested here? \\
\midrule
C1 & Theorem~8: HSP for hidden \emph{normal} subgroup of solvable
     black-box group is BQP.
   & indirectly (theoretical) & partly, via C4/C5 below \\
C2 & Theorem~11: HSP in $G$ is solvable in time
     $\mathrm{poly}(\text{input})+|G'|$.
   & indirectly (theoretical) & no \\
C3 & Theorem~13 statement: HSP in $G=N\rtimes (G/N)$ with $N$
     elementary abelian 2-group, $G/N$ cyclic, is
     poly-time BQP.
   & indirectly (theoretical) & \textbf{yes}, via C4/C5/C6 \\
C4 & \S6 reduction: the auxiliary $F$ hides a subgroup
     $K\le\mathbb{Z}_2\times N$.
   & \textbf{directly} & \textbf{yes} \\
C5 & Correctness of Fourier sampling on $\mathbb{Z}_2\times N$: all
     measurement outcomes lie in $K^\perp$.
   & \textbf{directly (exact)} & \textbf{yes} \\
C6 & Sample complexity: $O(\log|K|)$ samples suffice to recover $K$
     w.h.p.  ($K$ is generated by $O(\log|K|)$ random elements of
     $K^\perp$ dualised).
   & \textbf{directly} & \textbf{yes} \\
C7 & $|V|=O(\log|G/N|)$ coset representatives when $G/N$ is cyclic
     (\S6 paragraph after the theorem).
   & \textbf{directly} & spot-check, $G/N=\mathbb{Z}_2\Rightarrow|V|=1$ \\
\bottomrule
\end{tabular}
\end{center}

\section{Method}
\subsection{Group and hidden-subgroup construction}
\begin{enumerate}
  \item Fix $k\in\{2,3,4,5,6\}$ and $N=\mathbb{F}_2^k$ (elementary abelian
        2-group, $|N|=2^k$).
  \item Sample a random invertible $\sigma\in\mathrm{GL}_k(\mathbb{F}_2)$
        (three seeds per $k$; determinant check over $\mathbb{F}_2$).
  \item Compute a non-trivial fixed vector $u\ne 0$ of $\sigma$
        by finding a non-zero kernel element of $\sigma-I$ over
        $\mathbb{F}_2$ (necessary so that $H=\{(0,0),(u,1)\}$ is actually a
        subgroup: $(u,1)^2 = (u\oplus \sigma(u),0) = (0,0)$).
  \item Build $G=N\rtimes\mathbb{Z}_2$ with multiplication
        $(x_1,b_1)(x_2,b_2)=(x_1\oplus\sigma^{b_1}(x_2),\,b_1\oplus b_2)$.
  \item Choose the hidden subgroup $H=\{(0,0),(u,1)\}$ — the paradigmatic
        \emph{non-Abelian} case (an order-2 reflection subgroup not contained
        in~$N$), which is precisely the kind of subgroup the general
        Ettinger--H\o yer / IMS reduction is designed to catch.
  \item Define the HSP oracle $f\colon G\to \Sigma$ by assigning to each
        element the canonical label of its left coset $gH$
        (\texttt{coset\_label} in the code).
\end{enumerate}

\subsection{IMS \S6 reduction (auxiliary function $F$)}
The unique non-trivial coset representative is $z=(0,1)$ (since
$G/N=\mathbb{Z}_2$).  We build
\[
   F(0,x) = f((x,0)),\qquad F(1,x) = f((x,0)\cdot z) = f((x,1))
   \qquad (x\in\mathbb{F}_2^k).
\]
By direct enumeration one verifies that
\[
  F(b_1,x_1)=F(b_2,x_2)
  \iff x_1\oplus x_2 = \sigma^{b_1}(u)\cdot(b_1\oplus b_2)\,
  \bigl(\bmod 2\bigr),
\]
which defines the hidden subgroup $K=\langle (1,u)\rangle\le
\mathbb{Z}_2\times N$ of order~2 for every choice of $\sigma$ and $u$.

\subsection{Quantum statevector simulation}
Register layout: $(k{+}1)$ qubits.  We construct the reduced density matrix
$\rho$ on the group register that would arise from preparing
\[
   \frac{1}{\sqrt{2\cdot 2^k}}\sum_{b\in\mathbb{Z}_2,\,x\in N}
      |b\rangle|x\rangle|F(b,x)\rangle
\]
and tracing out the oracle-value register — explicitly $\rho = MM^\dagger$
where $M[i,v]=1/\sqrt{2^{k+1}}$ if $F$ evaluated at index $i$ equals value
$v$, else $0$.  We then apply
$\mathcal{F} = H^{\otimes (k+1)}$ (the QFT for $(\mathbb{Z}_2)^{k+1}$) and
read off the diagonal of $\mathcal{F}\rho\mathcal{F}^\dagger$ as the
measurement probability distribution over characters
$\chi_y$, $y\in\mathbb{F}_2^{k+1}$.

\subsection{Recovery procedure}
Given samples $y\in K^\perp$, we accumulate $y$-vectors into a
$\mathbb{F}_2$-basis and stop when its rank equals $\dim K^\perp = k{+}1-\log_2|K|$.
We then compute $K = (K^\perp)^\perp$ by brute enumeration
($2^{k+1}$ elements — small for $k\le 6$).

\subsection{Software \& versions}
\begin{itemize}
  \item Python 3.14.6 (Homebrew macOS)
  \item NumPy 2.5.1
  \item Qiskit 2.5.0
  \item Qiskit-Aer 0.17.2
  \item Custom code: \verb|code/ims_theorem13.py| (395 lines) and the earlier
        exploratory \verb|code/dihedral_hsp.py| (kept as an evidence artifact
        showing why plain $D_n$ with $n$ odd does \emph{not} fit \S6's
        assumptions — a useful negative control).
\end{itemize}

\subsection{Exact reproduction commands}
\begin{lstlisting}[basicstyle=\ttfamily\small,breaklines=true]
cd $HOME/Dropbox/REPLICATE-PROJECT/QC-200/QC-quant-ph-0102014-nonabelian-hidden-subgroup
python3 -m venv .venv && source .venv/bin/activate
pip install --quiet qiskit qiskit-aer numpy
python3 code/ims_theorem13.py     # writes code/ims_theorem13_results.json
python3 code/dihedral_hsp.py      # negative-control: plain D_n with n odd
\end{lstlisting}

\section{Results vs.\ paper}
Every one of the 15 test instances reproduces the exact behaviour
predicted by Theorem~13 / the standard Abelian-HSP correctness lemma.

\begin{center}\footnotesize
\begin{tabular}{@{}rrrrrrrl@{}}
\toprule
$k$ & seed & $|G|$ & $|K|$ & $|K^\perp|$ & $P[K^\perp]$ &
   $\max\!P[\text{off}]$ & samples used $\to$ match \\
\midrule
2 &  7 &   8 & 2 &  4 & 1.00000000 & 0.00e+00 & 2 $\to$ True\\
2 & 13 &   8 & 2 &  4 & 1.00000000 & 0.00e+00 & 4 $\to$ True\\
2 & 42 &   8 & 2 &  4 & 1.00000000 & 0.00e+00 & 2 $\to$ True\\
3 &  7 &  16 & 2 &  8 & 1.00000000 & 0.00e+00 & 3 $\to$ True\\
3 & 13 &  16 & 2 &  8 & 1.00000000 & 0.00e+00 & 6 $\to$ True\\
3 & 42 &  16 & 2 &  8 & 1.00000000 & 0.00e+00 & 6 $\to$ True\\
4 &  7 &  32 & 2 & 16 & 1.00000000 & 0.00e+00 & 4 $\to$ True\\
4 & 13 &  32 & 2 & 16 & 1.00000000 & 0.00e+00 & 6 $\to$ True\\
4 & 42 &  32 & 2 & 16 & 1.00000000 & 0.00e+00 & 6 $\to$ True\\
5 &  7 &  64 & 2 & 32 & 1.00000000 & 0.00e+00 & 5 $\to$ True\\
5 & 13 &  64 & 2 & 32 & 1.00000000 & 0.00e+00 & 9 $\to$ True\\
5 & 42 &  64 & 2 & 32 & 1.00000000 & 0.00e+00 & 8 $\to$ True\\
6 &  7 & 128 & 2 & 64 & 1.00000000 & 0.00e+00 & 8 $\to$ True\\
6 & 13 & 128 & 2 & 64 & 1.00000000 & 0.00e+00 & 9 $\to$ True\\
6 & 42 & 128 & 2 & 64 & 1.00000000 & 0.00e+00 & 8 $\to$ True\\
\bottomrule
\end{tabular}
\end{center}
Full JSON dump: \verb|report/evidence/ims_theorem13_results.json|.

\paragraph{Sample-complexity fit.}
Median samples-to-recover across seeds is
$(k{+}1)$ approx $\{3,5,6,7,8\}$ for $k=\{2,3,4,5,6\}$ — matching the
standard coupon-collector bound $\Theta(\log|K^\perp|)$ for random Fourier
sampling that Theorem~13 uses as a black box (via Theorem~3, Abelian HSP).

\paragraph{Negative control.}
\verb|dihedral_hsp.py| runs the same construction on the dihedral group
$D_n$ with $\mathbb{Z}_n$ replacing the elementary-abelian $N$.  For
$n\in\{3,5,7,\ldots\}$ (odd) or when $d\ne 0, n/2$, the linearisation
$K=\langle(1,d)\rangle$ is \emph{not} a subgroup of $\mathbb{Z}_2\times
\mathbb{Z}_n$ because $(1,d)+(1,d)=(0,2d\bmod n)\ne(0,0)$; hence the
Fourier-sampling probability mass ceases to sit on any exact lattice and
the linearised recovery attempt fails.  This is precisely why \S6
\emph{requires} $N$ to be an elementary abelian 2-group — a useful sanity
check that the reproducer respected the theorem's hypothesis.

\section{Discussion}
\begin{itemize}
  \item The reproduction confirms the \emph{correctness portion} of
        Theorem~13: the paper's reduction is not merely
        polynomial-time in the abstract, it is
        \emph{exact} at the level of measurement statistics.
  \item The paper's asymptotic claim (BQP) rests on the sample-complexity
        being $O(k)$ per non-trivial coset representative; our
        empirical median $(k{+}1)$ samples-to-full-recovery is consistent
        with this.
  \item We did not reproduce Theorem~11 (small commutator subgroup)
        because its efficient regime requires $|G'|$ small enough for
        classical enumeration but a group large enough to be non-trivial —
        a bigger engineering exercise than warranted for a QC-200 wave.
  \item The paper does not report any numerical benchmark, so this
        replication is inherently structural / correctness-oriented rather
        than a headline-number reproduction.
\end{itemize}

\section{Open Questions}
\begin{enumerate}
  \item[Q1.] The paper handles $G/N$ cyclic via Theorem~10 to
        produce a $\log|G/N|$-size $V$.  For non-cyclic $G/N$ they
        fall back to enumerating all $|G/N|$ coset reps.  Is there
        an intermediate regime (e.g.\ $G/N$ nilpotent of small class
        or metacyclic) in which a $\mathrm{poly}(\log|G/N|)$
        coset-rep set still exists but the paper's Theorem~13
        strategy is unable to exploit it?
  \item[Q2.] Our simulation always chose $H$ of order~2 (a single fixed
        reflection).  How does the sample-complexity curve
        (empirical median samples vs.\ $k$) behave when $H$ has larger
        order — i.e.\ $H=\{(0,0)\}\cup\bigl(u_1H_0\cup\cdots\cup u_r H_0\bigr)$
        with $H_0\le N$ of dimension $r\ge 1$?  The paper's bound
        $O(\log|G|)$ is worst-case; is the constant actually smaller
        for large-$H$ instances because $K$ is bigger and hence
        $|K^\perp|$ smaller?
  \item[Q3.] The reduction relies on the ability to \emph{compute}
        $f((x,0)\cdot z)$ from the same oracle. In a black-box quantum
        setting the query complexity of this composite $F$ is exactly
        twice that of $f$.  For a coherent-superposition oracle, is
        there a way to build $F$ with a single controlled-$z$ query
        rather than two disjoint calls, saving a factor of~2 in
        oracle queries?
  \item[Q4.] The \S6 argument crucially uses that $\mathbb{Z}_2\times N$
        is abelian only because $N$ is elementary abelian~2.  What is the
        smallest non-elementary-abelian $N$ (e.g.\ $N=\mathbb{Z}_4$ or
        $N=\mathbb{Z}_2\times\mathbb{Z}_4$) for which the same $F$
        construction still gives an abelian-HSP-solvable problem?  Our
        \verb|dihedral_hsp.py| negative control suggests the answer is
        \emph{never} for cyclic $\mathbb{Z}_n$ with $n$ odd — but is there
        a hybrid decomposition when $N$ is a mixed 2-group?
  \item[Q5.] The final sample-based recovery of $K$ is a classical
        post-processing step in our replication (Gaussian elimination
        over $\mathbb{F}_2$).  The paper is silent on whether a
        \emph{coherent} amplification (e.g.\ Grover-style amplitude
        amplification of the $K^\perp$-sampling stage combined with
        quantum linear-algebra) would provide any asymptotic improvement,
        or whether the abelian-HSP subroutine is already tight.
        Is there a quantum lower bound
        matching the $\Theta(\log|K|)$ upper bound achieved here?
\end{enumerate}
See \verb|report/open_questions.json| for machine-readable versions of
these five questions with per-question basis and follow-on next steps.

\section{Files in this replication}
See \verb|report/artifacts_summary.md| for the full inventory and
\verb|report/workflow.md| for a step-by-step reproduction guide.

\end{document}
